% SPDX-License-Identifier: Apache-2.0
%
% The flagship. Built by //docs:flagship. Every listing is included
% from the tree.
\documentclass[journal,9pt]{IEEEtran}

% Listings of Rust set by rustfmt at 80 columns: 5.6pt is what fits
% the frame of a column (issue 1061).
\newcommand{\listingsize}{\fontsize{5.6}{6.6}\selectfont}
% SPDX-License-Identifier: Apache-2.0
%
% The house preamble, shared by every document in this directory. Loaded
% after \documentclass, before \title. Keeping it in one file is what
% stops two articles drifting apart in font, listing style or figure
% style.
% Computer Modern. IEEEtran selects Times (ptm) by default, so the face has
% to be chosen explicitly.
%
% Latin Modern is the Computer Modern variety used here, and the choice is
% forced rather than aesthetic. Under T1 encoding, plain `cmr` has no Type 1
% outlines unless cm-super is installed, and it is not in the pinned
% distribution, so pdflatex falls back to EC bitmaps and every glyph in the
% PDF becomes a Type 3 font. That was measured: the first build produced a
% document with 10 Type 3 fonts and no Type 1. Latin Modern is Computer
% Modern redrawn with a full T1 repertoire and real .pfb outlines, so the
% same design comes out scalable.
\usepackage[T1]{fontenc}
\usepackage[utf8]{inputenc}
\usepackage{lmodern}
% microtype is not in the pinned TeX distribution. emergencystretch alone
% keeps the two-column measure from producing overfull lines.
\setlength{\emergencystretch}{3em}

\usepackage{amsmath}
\usepackage{amssymb}
\usepackage{amsthm}
\usepackage{booktabs}
\usepackage{array}
% enumitem is not in the pinned TeX distribution; the list spacing below
% does what \setlist would have done.
\usepackage{float}
\usepackage{xcolor}
\usepackage{listings}
\usepackage{tikz}
\usetikzlibrary{arrows.meta,positioning,fit,backgrounds,calc}
\usepackage[hidelinks,bookmarks=true]{hyperref}

% Numbered, definition-style box for a rule the reader is meant to apply.
\theoremstyle{definition}
\newtheorem{rulebox}{Rule}
\newtheorem{example}{Example}

% Unnumbered asides.
\theoremstyle{remark}
\newtheorem*{tip}{Tip}
\newtheorem*{pitfall}{Pitfall}

% Tighter lists, in place of enumitem's \setlist.
\setlength{\topsep}{2pt}
\setlength{\itemsep}{1pt}
\setlength{\parsep}{0pt}

% Code in the running text. The typewriter face under T1 ligates `<<`
% and `>>` into guillemets, so a shift printed as \code{<<} came out as
% a French quotation mark (issue 571). microtype's \DisableLigatures is
% not in the pinned distribution, but the primitive it rests on is
% pdfTeX's own: \pdfnoligatures strips every ligature from the font it
% is given, and the current font inside \texttt is the typewriter face
% at the size in use, so each size is stripped the first time \code
% sets it. Code wants no ligature at all: two quotes are two quotes.
\newcommand{\code}[1]{\texttt{\ifdefined\pdfnoligatures\pdfnoligatures\font\fi #1}}
\newcommand{\term}[1]{\emph{#1}}

% Syntax highlighting. Four roles, distinguishable in colour and still
% distinguishable in greyscale, because an IEEE article gets printed: the
% keyword is the only bold role, the comment is the only italic one, and the
% two remaining roles differ in lightness as well as in hue.
\definecolor{kwblue}{RGB}{20,60,150}
\definecolor{cmtgreen}{RGB}{90,120,90}
\definecolor{strred}{RGB}{150,50,40}
\definecolor{numgray}{gray}{0.45}
\definecolor{framegray}{gray}{0.70}
\definecolor{bgshade}{gray}{0.97}
% A darker fill for figure cells. The listing background has to stay very
% light to keep code readable; a figure cell has to be clearly shaded or the
% distinction it draws is invisible in print.
\definecolor{cellshade}{gray}{0.90}

% The listing size is the document's choice, because it depends on the
% column: a two-column article sets `\listingsize` to 6pt and keeps its
% listings within 70 columns, or to 5.6pt when it includes source that
% rustfmt sets at 80, which is what fits a column's frame (issue 1061);
% a one-column document keeps the body size and includes source files
% at 80. `//docs:frame_test` checks every listing line against its frame. Lines are never broken; a line that does not fit
% overflows the frame, which is visible, rather than wrapping, which is
% not.
\providecommand{\listingsize}{\scriptsize}

% `'` is deliberately not a string delimiter: the language uses it for loop
% labels, as Rust does, so treating it as a quote would swallow the rest of
% the line.
\lstdefinestyle{txhdl}{
  basicstyle=\ttfamily\listingsize,
  keywordstyle=\bfseries\color{kwblue},
  commentstyle=\itshape\color{cmtgreen},
  stringstyle=\color{strred},
  numberstyle=\tiny\color{numgray},
  morekeywords={mod,use,pub,const,type,struct,enum,trait,impl,fn,async,let,mut,
    match,if,else,loop,while,for,in,break,continue,return,as,await,
    module,bus,channel,transaction,protocol,pipeline,flow,seq,config,
    port,stage,pipe,instance,connect,bind,tag,domain,blackbox,foreign,
    property,role,signal,handler,true,false,
    sequence,interface,modport,implements,package,import,var,wire,func,proc,
    case,default,next,fork,branch,join,state,fsm},
  morecomment=[l]{//},
  morestring=[b]",
  showstringspaces=false,
  columns=fullflexible,
  keepspaces=true,
  breaklines=false,
  % Framed, so a listing is bounded rather than floating in the column.
  frame=single,
  rulecolor=\color{framegray},
  framesep=3pt,
  backgroundcolor=\color{bgshade},
  xleftmargin=3pt,
  xrightmargin=3pt,
  aboveskip=6pt,
  belowskip=6pt,
  % Every listing is numbered and captioned, so it can be referred to by
  % number. `caption` without `float` numbers the listing and leaves it
  % where it was written, which is what a listing in running prose needs.
  captionpos=b,
  abovecaptionskip=2pt,
  belowcaptionskip=6pt,
}

% Figures drawn as text. Framed the same way, and with the highlighting
% switched off, because a box-and-arrow diagram is not code and half its
% words turning blue reads as an accident. Setting a style adds keys rather
% than clearing the ones already set, so the three roles are neutralised by
% hand rather than by leaving them out. Program output is printed in this
% style too, and a netlist's assignment can run past the frame, so a line
% that does not fit is broken and the rest indented; source never is.
\lstdefinestyle{ascii}{
  basicstyle=\ttfamily\listingsize,
  keywordstyle=\ttfamily,
  commentstyle=\ttfamily,
  stringstyle=\ttfamily,
  columns=fullflexible,
  keepspaces=true,
  breaklines=true,
  breakindent=2em,
  frame=single,
  rulecolor=\color{framegray},
  framesep=3pt,
  backgroundcolor=\color{bgshade},
  xleftmargin=3pt,
  xrightmargin=3pt,
  aboveskip=6pt,
  belowskip=6pt,
}
\lstset{style=txhdl}

% House TikZ styles. `shorten` lives on the arrow styles rather than on each
% arrow, so every arrow in the document gets the gap, including ones added
% later. TikZ clips a path at a node's border, so without it an arrowhead
% butts against the box with no gap at all. Keep `node distance` well above
% twice the shorten, or the arrow shrinks to nothing.
\tikzset{
  box/.style={draw,rounded corners=2pt,align=center,inner sep=4pt,
              font=\footnotesize},
  boxf/.style={box,fill=bgshade},
  lbl/.style={font=\scriptsize\itshape,align=center},
  fl/.style={-{Stealth[length=4pt]},shorten >=2.5pt,shorten <=2.5pt},
  fld/.style={fl,dashed},
}


% The contents list: the class gives a section's number 2.75em, which
% a Roman numeral past thirty overruns onto the title, so the box is
% wider here, and a subsection's entry starts where a title does.
\makeatletter
\def\l@section#1#2{\addpenalty{\@secpenalty}\addvspace{1.0em plus 1pt}%
    \@tempdima 5.75em \begingroup \parindent \z@ \rightskip \@pnumwidth%
    \parfillskip-\@pnumwidth {\bfseries\leavevmode #1}\hfil\hbox to\@pnumwidth{\hss #2}\par%
    \endgroup}
\def\l@subsection{\@dottedtocline{2}{5.75em}{5em}}
\def\l@subsubsection{\@dottedtocline{3}{10.75em}{5.5em}}
\makeatother

% SPDX-License-Identifier: Apache-2.0
% The boards' memory maps, written by //tools/memmap from the maps the
% designs decode (issue 444). A document asks for a table with
% \mmget{<what>}{<board>}, and a table the generator does not write
% stops the document rather than leaving a gap.
\input{tools/memmap/mm_tables}
\makeatletter
\newcommand{\mmget}[2]{%
  \@ifundefined{mm@#1@#2}%
    {\PackageError{memmap}{//tools/memmap writes no #1 for #2}%
      {Add it to tools/memmap/main.rs.}}%
    {\csname mm@#1@#2\endcsname}}
\makeatother


% The colours of the placement map, one per subsystem, chosen to be
% told apart in print as well as on a screen. //tools/fpgamap names
% them, so they are defined here and nowhere else.
\definecolor{flagdie}{gray}{0.92}
\definecolor{flagcore}{RGB}{40,90,170}
\definecolor{flagmem}{RGB}{200,110,30}
\definecolor{flageth}{RGB}{50,130,60}
\definecolor{flagvid}{RGB}{140,60,150}
\definecolor{flagcdc}{RGB}{30,150,160}
\definecolor{flagtop}{gray}{0.45}

\InputIfFileExists{buildstamp.tex}{}{}
\providecommand{\buildstamp}{Draft build (outside the Bazel flow).}

\title{The TxHDL Flagship}

\author{Filip Filmar and Dragiša Janković%
\thanks{Both authors are independent. Filip Filmar is the
corresponding author, at f@hdlfactory.com; Dragiša Janković is
reached at linkedin.com/in/dragisa-jankovic.}%
\thanks{This document describes the unified hardware platform and peripheral
subsystems~\cite{eth,hdmi,axi} running on the Vreteno processor~\cite{vreteno}.
All listings are embedded from the project repository \code{HDL/txhdl} during
compilation. The project overview~\cite{cover} indexes all companion reports.}%
\thanks{The authors used Claude, an AI assistant, while developing
architectural concepts, documentation, and software. Verbatim prompts are
preserved in repository commit logs.}%
\thanks{\buildstamp}}

\markboth{The TxHDL Flagship}{Filmar and Janković: The TxHDL Flagship}

\begin{document}
\maketitle

\begin{abstract}
Hardware subsystems in this project were initially verified on the Alinx
AX7A200B FPGA board as standalone designs: a processor with external DDR3
memory, an Ethernet echo transceiver, and an HDMI video generator. Each
experiment required a dedicated bitstream. The flagship combines these
independent components into a unified bitstream written to on-board QSPI
flash, establishing a persistent baseline upon board power-up. The design also
stabilizes the hardware foundation. The processor boots into a ROM loader that
accepts software binaries over the serial interface in seconds, leaving
underlying gateware undisturbed. This document analyzes subsystem
integration, clock domain crossing topologies, diagnostic status indicators,
and non-volatile flash deployment.
\end{abstract}

\section{Why a flagship}
\label{sec:f-why}

An FPGA device holds one configuration bitstream at a time. Evaluating
peripherals through standalone demonstrations left the physical board in an
ephemeral state between test runs. Demonstrating the complete platform
required reconfiguring the FPGA repeatedly for each subsystem.

The flagship resolves this fragmentation. It consolidates all verified
subsystems into a single bitstream stored in non-volatile QSPI flash,
rebuilding whenever constituent logic changes. Applying power restores the
complete platform immediately without requiring JTAG initialization.

This architecture also accelerates software iteration. Full FPGA placement
and routing for the combined system requires thirty minutes. Embedding software
images directly into synthesis netlists would impose a thirty-minute turnaround
on every code change. Instead, the processor boots into a compact ROM
loader~\cite{vreteno}. The loader listens on the serial port, receives an
executable binary, writes it to DDR3 memory, and branches to it. Because the
hardware netlist remains fixed, updating software takes seconds.

\section{What is in it}
\label{sec:f-what}

The flagship integrates three verified subsystems without altering their
internal logic.

\emph{Processor and memory.} The \code{board} netlist emitted by
\code{\#[lower]} for \code{vreteno32::board::Board} provides a Vreteno RV32IMAC
processor, an AXI bus tracker, an interconnect router, data memory, machine
timers, software interrupts, a platform-level interrupt controller, a
pulse-width modulator, a 16550-compatible UART, and one gigabyte of DDR3
DRAM managed by an AMD MIG controller. Processor boot ROM contains the serial
loader.

\emph{Ethernet interface.} Lowered MAC transmit and receive units connect to
external physical layer transceivers through an RGMII wrapper~\cite{eth}. Behind
the MAC sits the remote peripheral mapped at \code{0x3300}, which proxies bus
accesses to a host workstation. A transaction issued by the processor departs
within an Ethernet frame and awaits a matching response frame. This topology
allows developers to evaluate peripheral drivers in host software against
authentic bus semantics before writing dedicated RTL. Section~\ref{sec:f-remote}
details this interface. The processor also accesses the MAC directly through
\code{EthSlots} at peripheral slot five, transmitting and receiving raw network
frames. Hardware validation verified outbound frame counts on the link server,
and \code{ethrx} and \code{ethirq} passed regression tests. Furthermore, a
Zephyr fastboot service transferred over the serial link operated on UDP port
5554, responding to \code{getvar} queries and booting target images via
\code{fastboot boot}.

\emph{HDMI display output.} The video pipeline pairs a display framebuffer
with an I2C controller that initializes the on-board SiI9134 HDMI
transmitter~\cite{hdmi}. The flagship couples this pipeline to the system
bus, allowing software to update display surfaces and drive external monitors.

\section{The third slot}
\label{sec:f-slot}

The system router decodes eight master ports, listed in Table~\ref{tab:f-map}.
Address zero maps read-only boot ROM, leaving no spare router ports. Compact
peripherals multiplex the 4~KiB page at \code{0x3000} through an AXI-Lite
bridge, each occupying a 256-byte sub-slot. This shared mapping prevents
low-density register banks from exhausting router ports. The video peripheral
exposes three control registers and occupies slot three at \code{0x3200}, as
detailed in Table~\ref{tab:f-page} and Table~\ref{tab:f-hdmi}. Tooling at
\code{//tools/memmap} generates these tables directly from router and decoder
definitions, matching the synthesized hardware map.

\begin{table}[htbp]
\caption{The flagship's address map, lowest first, as the core's
router decodes it.}
\label{tab:f-map}
\mmget{board}{flagship}
\end{table}

\begin{table}[htbp]
\caption{The peripheral page at \code{0x3000}, as its bridge decodes
it.}
\label{tab:f-page}
\mmget{page}{flagship}
\end{table}

\begin{table*}[tbp]
\caption{The video peripheral's registers, on the third slot.}
\label{tab:f-hdmi}
\mmget{regs}{hdmi}
\end{table*}

What is different about the third slot is that it does not reach a
field of \code{Board}. It leaves the unit as five channel ports:

\lstinputlisting[rangebeginprefix=//\ begin\{,rangebeginsuffix=\},%
rangeendprefix=//\ end\{,rangeendsuffix=\},includerangemarker=false,%
linerange=vslot-vslot,caption={\code{cpu/vreteno/src/board.rs}: the
third slot, and why it leaves the
unit.},label={lst:f-slot}]{cpu/vreteno/src/board.rs}

This external routing accommodates an independent clock domain. The video
scanout logic requires a 25.2\,MHz pixel clock, whereas the processor operates
at 100\,MHz. Placing the video controller inside \code{Board} would constrain
it to the processor clock, violating standard VGA timing. Exporting the slot
allows the top-level module to bridge clock boundaries. Target configurations
lacking video hardware tie these channel ports off, as implemented in
\code{cpu/vreteno/board/vreteno\_board.v}.

Slot four provides an instructive contrast. The remote peripheral operates
synchronously on the 100\,MHz processor clock, so it resides directly within
\code{Board} alongside the serial controller. Rather than exporting the bus
slot, the unit exports the peripheral's external interface: packet data and
framing delimiters. Subsystems export bus slots only when downstream hardware
requires an independent clock.

\section{Three clocks, one input}
\label{sec:f-clocks}

The carrier board supplies a single 200\,MHz differential clock input. The
flagship buffers this reference once and feeds three dedicated clock generators,
matching the configurations verified in the individual subsystem designs:

\begin{itemize}
\item The DDR3 memory controller clock generator uses the 200\,MHz input as its
  delay reference, runs internal DRAM logic at 400\,MHz, and generates a
  100\,MHz synchronous system clock for processor execution and the serial
  interface.
\item An MMCM multiplies the input by five to 1000\,MHz and divides by eight to
  produce a 125\,MHz clock for the Ethernet transmit path. An internal MMCM in
  \code{eth\_rgmii} shifts the recovered receive clock by 78.75 degrees
  (1.75\,ns), centering sampling edges within data windows presented by the
  physical layer transceiver. Empirical measurements indicate the PHY
  introduces approximately 2\,ns of internal receive delay.
\item A video MMCM scales the reference clock to 25.2\,MHz for pixel scanout,
  matching the standard 640$\times$480 60\,Hz display rate within 0.1\% of
  nominal (25.175\,MHz).
\end{itemize}

Two data paths cross clock boundaries, both implemented as FIFO channels rather
than unsynchronized scalar signals. The first path routes the processor's
AXI-Lite interface to the video peripheral. Each of the five AXI-Lite channels
crosses independently through \code{chan\_cdc}, an asynchronous FIFO with
Gray-coded read and write pointers. Routing each channel independently satisfies
all protocol requirements because AXI channels maintain no inter-channel cycle
dependencies; a write transaction completes when its write response handshake
returns.

The second path transfers remote peripheral packets between clock domains,
routing bytes and end-of-frame flags across the boundaries. Outbound packets
cross from the 100\,MHz processor clock to the 125\,MHz transmit clock, and
inbound packets cross from the recovered PHY clock to the processor clock. The
inbound FIFO allocates 128 bytes rather than the 16 bytes used in earlier echo
tests. Because the Ethernet MAC delivers one byte every 8\,ns while the
processor reads one byte every 10\,ns, the FIFO absorbs packet bursts while the
software reader catches up.

Listing~\ref{lst:f-cdc} presents \code{lib/board/chan\_cdc.v}, which implements
the dual-clock channel synchronizer. Common board infrastructure now resides
under \code{//lib/board} for reuse across top-level designs.

\lstinputlisting[caption={\code{lib/board/chan\_cdc.v}: a lowered
unit's channel, transferred across clock
domains.},label={lst:f-cdc}]{lib/board/chan_cdc.v}

\section{The peripheral that is a program}
\label{sec:f-remote}

Developing custom peripheral hardware requires extensive simulation, synthesis,
placement, and routing before physical register behaviors can be validated. The
remote peripheral architecture decouples register definition from FPGA
compilation. \code{Remote} occupies peripheral slot four, behaving like any
standard memory-mapped slave while handling no register operations locally.
Accepted bus transactions dispatch over an outbound channel, awaiting
completions on a matching response channel. \code{RemoteLink} marshals outgoing
requests into 27-byte Ethernet frames and demarshals inbound replies, passing
them through the MAC to the network. The processor communicates directly with a
software service hosted on an attached workstation.

\lstinputlisting[rangebeginprefix=//\ begin\{,rangebeginsuffix=\},%
rangeendprefix=//\ end\{,rangeendsuffix=\},includerangemarker=false,%
linerange=remote-remote,caption={\code{cpu/vreteno/src/board.rs}: the
peripheral and the link that puts it on the
port.},label={lst:f-remote}]{cpu/vreteno/src/board.rs}

Two parameters are fixed within the RTL architecture. The device identifier
is 1: every request frame holds this tag, allowing target software to multiplex
multiple devices on a single physical link. Transaction patience is set to one
second (100\,000\,000 cycles at 100\,MHz). If no response frame arrives before
the timer expires, the peripheral asserts an AXI slave error (\code{SlvErr})
to prevent bus stalls.

The host software endpoint resides in \code{//tools/remote}, a Go service that
processes incoming frames, implements peripheral register behaviors, and
returns response packets. The service exchanges length-prefixed packet streams
across an ssh tunnel. On the target-attached host, \code{//tools/remote/shim}
captures peripheral frames via raw Linux network sockets, encapsulates them
with length headers into the stream, and dispatches return frames onto the
physical link.

Full system integration succeeded on September 22, 2026. With the flagship
bitstream loaded into flash, the bridge service running on the target host, and
\code{//tools/remote} registered as device~1, the loader dispatched the
diagnostic binary in Listing~\ref{lst:f-remote-rs}. The processor executed two
register stores, read back both values over Ethernet, and confirmed data
integrity on the serial console. Listing~\ref{lst:f-remote-run} records the
console output and server transaction logs.

Testing exposed two timing constraints. Initial gateware used a 20\,ms response
timeout (\code{REMOTE\_WAIT}). When accessing development services across an ssh
tunnel with a 175\,ms round-trip latency, transactions timed out prematurely,
returning zero. Extending timeout patience to one second (issue~\#397) restored
robust operation across network tunnels. Additionally, serial warm resets
originally failed to reinitialize processor halt registers; clearing these
registers upon reset (issue~\#398) ensures repeatable warm restarts.

\lstinputlisting[caption={\code{cpu/vreteno/rust/remote.rs}: two
words to the peripheral and back, with the verdict on the serial
port.},label={lst:f-remote-rs}]{cpu/vreteno/rust/remote.rs}

\lstinputlisting[style=ascii,caption={\code{docs/remote\_run.txt}:
what the board said, and what the program beside the shim saw, on
September 22, 2026.},label={lst:f-remote-run}]{remote_run.txt}

\section{The lights}
\label{sec:f-leds}

The carrier board features four active-low LEDs. The flagship assigns them to
critical status signals to facilitate hardware diagnostics:

\begin{itemize}
\item \textbf{LED1}, pulse-width modulator channel zero or core halt status. This
  indicator serves software applications: the boot loader leaves it unlit,
  allowing running software to control brightness or signal completion.
\item \textbf{LED2}, Ethernet frame activity. Pulse-stretching logic asserts
  this LED for 200\,ms upon detecting inbound or outbound packet traffic,
  ensuring visual visibility.
\item \textbf{LED3}, DDR3 calibration status. The MIG controller asserts this
  line upon completing memory calibration. Because the boot loader stages
  incoming software into DDR3 RAM, an unlit LED3 immediately isolates memory
  initialization failures.
\item \textbf{LED4}, system heartbeat. Toggling at 3\,Hz from the memory
  controller clock, this indicator verifies clock generation and active gateware
  execution.
\end{itemize}

The HDMI video pipeline omits dedicated diagnostic LEDs, as output monitors
display video status directly.

\section{The design}
\label{sec:f-top}

The top-level integration module brings together I/O pins, clock generators,
reset distribution networks, clock domain crossing FIFOs, and status
indicators.

\lstinputlisting[caption={\code{flagship/board/flagship.v}: the
flagship's top.},label={lst:f-top}]{flagship/board/flagship.v}

Physical constraints are shared across modules to eliminate duplication. Clock,
reset, LED, UART, and DDR3 pin assignments derive from constraint files under
\code{cpu/vreteno/board/}. Physical layer Ethernet pins import from
\code{eth/board/eth\_pins.xdc}, and HDMI encoder pins import from
\code{hdmi/board/hdmi\_pins.xdc}. Each FPGA pin appears in exactly one primary
constraint file, ensuring that pin modifications propagate consistently across
all build configurations. The flagship's dedicated constraint file defines only
inter-clock timing relationships, declaring asynchronous clock groups across
independent generators.

\lstinputlisting[caption={\code{flagship/board/ax7a200\_flagship.xdc}:
what no subsystem owns.},label={lst:f-xdc}]%
{flagship/board/ax7a200_flagship.xdc}

\section{Building it, and putting it in the flash}
\label{sec:f-build}

Vivado targets in this repository are marked manual to prevent unintended
long-running actions during routine builds. Compiling
\code{//flagship:flagship\_synth} requires five minutes, and subsequent place
and route via \code{//flagship:flagship\_pnr} consumes twenty-five minutes.
Because the implementation step depends directly on synthesis output, executing
place and route drives the complete flow.

Synthesis stops on warnings that would otherwise escape to hardware. Vivado
truncates sized literals exceeding declared bit widths and continues without an
error: \code{21'd2\_100\_000} in earlier top logic was truncated to 2848,
preventing the downstream reset counter from asserting. The build completed
without timing violations despite generating a defective bitstream, detectable
only through an obscure line in \code{flagship\_synth.log}. Detecting the failure
required a complete place-and-route cycle and board bring-up. Every synthesis
rule now queries Vivado after \code{synth\_design} to count occurrences of
critical warning codes, failing the build if counts are non-zero. These warning
rules are centralized in \code{tools/vivado/defs.bzl}.

Programming the board requires a TCP tunnel to the host attached to the FPGA,
opened via \code{bazel run //cpu/vreteno/board/remote:hw\_server}. Target
\code{//flagship:flagship\_prog} downloads bitstreams directly to volatile
FPGA SRAM for temporary testing, whereas \code{//flagship:flagship\_flash}
writes non-volatile QSPI flash for persistent deployment. Both targets accept
\code{-{}-hostport localhost:3122} and a hardware target filter.

Executing software returns to the loader through three coordinated mechanisms:
The carrier RESET button and adjacent KEY1 trigger core resets. External memory
controllers continue refreshing DRAM through the reset assertion, allowing
in-flight bus transactions to drain gracefully and preserving RAM contents
across warm restarts. An earlier revision reset the memory controller
simultaneously, stranding the AXI bridge and freezing subsequent boot attempts.
From the host workstation, pulling the serial RX line low for 30\,ms initiates
a warm reset. Invoking \code{load -{}-reset} transmits this low-speed pulse at
300 baud, re-entering the boot loader prior to transferring executable
binaries. Standard serial break conditions are avoided because the on-board
CP2102N USB-UART bridge drops host break pulses under Linux.
Core resets simultaneously clear MAC state machines and video logic through
two-stage synchronizing registers. Cross-domain FIFOs reset synchronously: each
\code{chan\_cdc} instance holds itself empty while either domain asserts reset,
discarding in-flight words rather than exposing stale data to reset logic
(issues 509 and 1325).

The non-volatile storage chip is a Micron MT25QL128 quad-SPI flash. The bitstream
occupies flash addresses beginning at zero, followed by software images at
offset \code{0x00A0\_0000} (issue~312). For this FPGA device, uncompressed
configuration bitstreams require 77\,845\,216 bits (9\,730\,652 bytes), ending
at address \code{0x0094\_7A5C}. Software partitions align with the subsequent
64~KiB erase sector at 10~MiB, leaving 6~MiB of flash capacity. The processor
accesses flash memory through an AXI window at \code{0x20A0\_0000}. Automated
tests \code{//flagship:flagship\_fits\_test} and
\code{//cpu/vreteno:vreteno\_board\_fits\_test} inspect bitstream headers to
confirm that configuration data does not exceed reserved partitions.
Furthermore, \code{//flagship:flagship\_flash} embeds the diagnostic program
\code{flashid} via \code{write\_cfgmem -loaddata}, validated by
\code{//flagship:flagship\_image\_test}.

With the flagship flashed, software deployment requires no FPGA compilation:
developers build binaries via \code{bazel build
//cpu/vreteno/rust:hello\_ram\_bin} and transmit them over the serial link
using \code{load}. This workflow reduces turnaround cycles from thirty minutes
to seconds.

\section{What it costs}
\label{sec:f-cost}

The complete flagship consumes a small fraction of the Xilinx XC7A200T FPGA.
Post-synthesis reports for git revision \code{806a229} record 10\,999 slice LUTs
(8.2\% utilization), 8\,918 slice registers (3.3\%), 16.5 dual-port 36~Kb Block
RAM tiles (4.5\% of 365), 4 DSP48E1 slices, and 126 bonded I/O pins (44\% of
285). Clock resources comprise one PLL and four MMCM instances, including the
dedicated MIG memory clocking structures. Logic capacity is not a design
bottleneck; subsequent subsystem integration depends strictly on closing timing
at target frequencies rather than conserving slices.

\section{Where it landed}
\label{sec:f-place}

Resource counts quantify overall logic consumption without capturing spatial
distribution. Multi-subsystem integration requires examining whether logic forms
cohesive functional regions or scatters across the die, and how pin assignments
influence placement.

Figure~\ref{fig:f-map} depicts cell placement extracted directly from routed
design checkpoints. Rule \code{//flagship:place\_report} attributes placed logic
cells to subsystem hierarchies, recording coordinates in
\code{docs/flagship\_place.tsv}. Build tool \code{//tools/fpgamap} rasterizes
these totals into the placement diagram. Table~\ref{tab:f-place} tabulates the
corresponding cell counts.

\begin{figure}[t]
\centering
\input{flagshipmap}
\caption{Where the flagship's subsystems landed on the xc7a200t. A
square is eight slices by eight; its colour is whichever subsystem
holds the most cells in it and its shade how full it is. The grey is
part that holds nothing.}
\label{fig:f-map}
\end{figure}

\begin{table}[t]
\caption{The same placement as numbers. A slice holds several cells,
so these counts are larger than a utilisation report's, which counts
slices.}
\label{tab:f-place}
\centering
\footnotesize
\input{flagshiptotals}
\end{table}

The placement map highlights four structural characteristics:

\emph{Subsystems form localized regions.} The memory controller occupies columns
133 to 163 along the right die boundary adjacent to memory I/O banks. The
processor core occupies central columns 97 to 145. Video logic forms compact
clusters on the left margin near HDMI transmitter pins. Cell placement
naturally aligns with external I/O interfaces and interconnect targets.

\emph{Ethernet logic migrates toward bus masters.} In standalone echo designs,
MAC logic clustered at columns 0 to 22 beside the PHY interface. Connected to
the system bus, the MAC shifts inward to columns 36 to 104, reflecting the
routing demands of its bus interconnect.

\emph{The memory controller dominates peripheral resources.} The controller
consumed 9\,493 cells compared to 10\,719 cells for the processor and
peripherals combined under UberDDR3, remaining locked against DRAM I/O banks.

\emph{Clock domain crossings cluster along subsystem boundaries.} The five FIFOs
transferring processor AXI-Lite transactions to the pixel clock occupy columns
104 to 122 between the processor and video pipelines. Frame FIFOs sit along the
Ethernet boundary nearest the core.

The map and the table reflect git revision \code{806a229}, routed from the same
checkpoint evaluated above.

\section{The first program down the wire}
\label{sec:f-first}

Display interfaces require visual validation. The initial demonstration software
renders a rotating icosahedron with the TxHDL logo in the corner
(Listing~\ref{lst:f-ico}), originally rendered directly by the processor and
subsequently accelerated by Razboj. Compiling for DDR3 memory at
\code{0x4000\_0000} and uploading via the serial loader enabled rapid graphics
iteration without repeated FPGA bitstream synthesis.

The rendering implementation features three key techniques. Rather than storing
a static face table, initialization routines identify all vertex triples
separated by pairwise unit edges, orienting face normals outward. Face normals
are unit vectors transformed by object rotation, allowing a single dot product
to evaluate visibility and Lambertian illumination. For convex polyhedra,
backface culling resolves visibility without a depth buffer. Early
processor-rendered versions composed display rows in a local scanline buffer
before writing to the framebuffer, avoiding visible flicker caused by drawing
directly into an active scanout buffer.

Visual evaluation on physical monitors revealed two rendering subtleties.
Positioning the virtual camera three radii away caused nearest faces to dominate
the silhouette, creating unnatural perspective distortion. Additionally,
independent color channel quantization introduced hue shifts during shading
transitions. The shader was revised to use a perceptual square-root intensity
ramp with an elevated black floor and $2\times2$ ordered dithering to smooth
four-bit color transitions.

Hardware rasterizer Razboj now accelerates rendering at $640\times480$ with
24-bit color depth (issue~\#986). The processor manages transformation matrices
and emits display lists (Listing~\ref{lst:f-ico-list}) containing background
clear rectangles and shaded triangles. Polygon winding order determines
screen-space visibility. The display pipeline utilizes alternating framebuffers
in memory. Razboj renders to the off-screen surface; upon completion, software
redirects scanout pointers during vertical blanking intervals to eliminate
display tearing. The processor renders the corner logo into both framebuffers
during initialization. Every 64 frames, software logs cycle counts for display
list generation, rasterization, and total frame latency.

\lstinputlisting[caption={\code{cpu/vreteno/rust/ico\_hdmi.rs}: the
icosahedron.},label={lst:f-ico}]{cpu/vreteno/rust/ico_hdmi.rs}

\lstinputlisting[caption={\code{cpu/vreteno/rust/ico\_list.rs}: the
icosahedron as Razboj's display lists.},label={lst:f-ico-list}]{cpu/vreteno/rust/ico_list.rs}

The logo is stored as pre-rasterized data in a dedicated crate
(Listing~\ref{lst:f-logo}). Sized at $144\times144$ pixels, it maps one-to-one
to screen pixels in the corner of the $640\times480$ display (issue~\#1213).
Pixel values are 4-bit indices into a 16-color palette generated by k-means
clustering, reducing image storage to 10~KB. Conversion utility
\code{//tools/logo2rs} (Listing~\ref{lst:f-logo2rs}) downsamples source images
and generates Rust source tables. Index zero represents transparent pixels,
preserving background colors underneath. Tables are declared as static arrays
rather than constants; in unoptimized builds, indexing constant arrays triggers
full array copies on every pixel access, inflating paint latency from two
seconds to 31 seconds (issue~\#1245).

\lstinputlisting[caption={\code{lib/logo/logo.rs}: the head of the
logo, as \code{logo2rs} writes it.},label={lst:f-logo},%
linerange={1-51}]{lib/logo/logo.rs}

\lstinputlisting[caption={\code{tools/logo2rs/main.go}: the PNG to the
crate.},label={lst:f-logo2rs},tabsize=2]{tools/logo2rs/main.go}

\section{What is not done}
\label{sec:f-todo}

\emph{Display pipeline maturity.} Section~\ref{sec:f-first} demonstrates
single-buffered scanout with manual tearing prevention and convex geometry
rendering without hardware depth buffering. Future graphics workloads will
require hardware double buffering and a depth buffer.

\emph{Target utilization headroom.} Architectural plans allow partitioning the
flagship into numbered configurations (\code{flagship-0}, \code{flagship-1}) if
logic density saturates the FPGA. As Section~\ref{sec:f-cost} documents, existing
logic consumes less than 10\% of device resources, keeping a unified flagship
viable.

\begin{thebibliography}{9}

\bibitem{vreteno}
F.~Filmar and D.~Janković, ``Vreteno in TxHDL,'' project repository
\code{HDL/txhdl}, \code{//docs:vreteno}, 2026.

\bibitem{eth}
F.~Filmar and D.~Janković, ``Ethernet in TxHDL,'' project repository
\code{HDL/txhdl}, \code{//docs:eth}, 2026.

\bibitem{hdmi}
F.~Filmar and D.~Janković, ``HDMI in TxHDL,'' project repository
\code{HDL/txhdl}, \code{//docs:hdmi}, 2026.

\bibitem{axi}
F.~Filmar and D.~Janković, ``AXI in TxHDL,'' project repository
\code{HDL/txhdl}, \code{//docs:axi}, 2026.

\bibitem{cover}
F.~Filmar and D.~Janković, ``TxHDL documents,'' project repository
\code{HDL/txhdl}, \code{//docs:cover}, 2026.

\end{thebibliography}

\end{document}
